\section{Cycle cells, circuit partitions, and an interlace representation}
\label{sec:random-cycle-quotient}

This section records an exact representation and a research direction.
It does not prove a new upper bound for the remaining graphs.
A chosen two-factor is required.
Its cycles need not be shortest cycles, and shortest cycles need not form a two-factor.

Let $G$ be a connected cubic graph with a specified two-factor consisting of $m$ cycles, called cells.
The edges outside the two-factor form a perfect matching.
Contract each cell to a vertex, retaining its cyclic order of matching incidences.
This produces a connected quotient multigraph $Q$, possibly with loops.
For an even edge set $F$ of $Q$, let $V_+(F)$ denote its nonisolated vertices and let $z(F)=m-|V_+(F)|$.
At a vertex of selected degree $2q>0$, the incident selected half-edges inherit a cyclic order.
Their two alternating pairings are
\[
 (12)(34)\cdots(2q-1,2q),
 \qquad
 (23)(45)\cdots(2q,1).
\]
When $q=1$ these pairings coincide, but the two cell paths are distinct and must both be counted.

For a choice $\sigma\in\{0,1\}^{V_+(F)}$ of alternating pairings, let
$\ell(F,\sigma)$ be the number of circuits after smoothing every selected quotient vertex according to its chosen pairing.
Define the component polynomial of the even space by
\[
 Z_G(t)=\sum_{J\text{ even in }G}t^{\operatorname{comp}_+(J)},
\]
where only components with edges are counted and the empty state has exponent zero.

\begin{proposition}[The exact cycle-cell expansion]
\label{prop:random-cycle-cell-expansion}
One has
\[
 Z_G(t)=
 \sum_{F\text{ even in }Q}
 (1+t)^{z(F)}
 \sum_{\sigma\in\{0,1\}^{V_+(F)}}t^{\ell(F,\sigma)}.
\]
For $F=\varnothing$, the inner sum is interpreted as one.
In particular,
\[
 \beta(G)=\beta(Q)+m,
 \qquad
 c(G)=m+
 \sum_{\substack{F\ne\varnothing\text{ even in }Q\\
                 (V_+(F),F)\text{ connected}}}\kappa(F),
\]
where $\kappa(F)$ counts the alternating-pairing choices producing one circuit.
\end{proposition}

\begin{proof}
The matching edges of an even state in $G$ induce an even boundary set at each cell, hence an even edge set $F$ of $Q$.
Conversely, any even boundary pattern on a cycle has exactly two extensions using cycle edges.
For a nonzero pattern they are complementary unions of paths, inducing the two alternating pairings.
For the zero pattern they are the empty set and the entire cell cycle.
All cell choices are independent.

Thus every even quotient state has exactly $2^m$ lifts, and its component polynomial is the displayed summand.
This also proves the rank identity.
A lift with one component either is one cell cycle and has $F=\varnothing$, or has nonempty connected quotient support.
In the latter case all zero-boundary cells must choose the empty extension, and the selected smoothing must have one circuit.
This gives the cycle-count identity.
\end{proof}

\begin{proposition}[A matrix on blocks of switches]
\label{prop:random-cycle-cell-interlace}
Fix a connected nonempty even quotient state $F$.
Suppose $\kappa(F)>0$ and choose one permitted one-circuit state as reference.
There is a symmetric matrix $M_F$ over $\mathbb F_2$ whose index set is partitioned into blocks $B_v$, one for each $v\in V_+(F)$, with
\[
 |B_v|=\frac{\deg_F(v)}2-1,
 \qquad
 \dim M_F=|F|-|V_+(F)|=\beta(F)-1,
\]
such that
\[
 \kappa(F)=
 \#\left\{
 S\subseteq V_+(F):
 M_F\!\left[\bigcup_{v\in S}B_v\right]\text{ is nonsingular}
 \right\}.
\]
Empty blocks remain in the subset count.
The empty principal matrix is nonsingular by convention.
\end{proposition}

\begin{proof}
Replace each selected vertex of degree $2q\ge4$ by a chain of $q-1$ four-valent switches.
This replacement has two permitted states, obtained by choosing one local transition at every switch in the chain or the other transition at every switch.
To verify the replacement, split off ports $1,2,3$ at the first switch, whose fourth incidence leads to the remaining gadget.
The first transition pairs $1$ with $2$ and sends $3$ into the remaining gadget.
The second pairs $2$ with $3$ and sends $1$ into it.
The remaining cyclic port list is the new internal incidence followed by $4,5,\ldots,2q$.
Induction gives the two alternating pairings.
For $q=1$, suppress the vertex and retain its twofold multiplicity.

The resulting graph is connected and four-regular unless every selected degree is two.
The reference state gives an Euler circuit.
Form its binary interlace matrix.
At each switch, put a one on the diagonal exactly when its alternative transition is orientation-inconsistent with that Euler circuit.
Call the resulting symmetric matrix $M_F$.

Changing a cell's state changes every switch in its block.
The extended Cohn--Lempel equality states that the new circuit count is
\[
 1+\operatorname{nullity}_{\mathbb F_2}
 M_F\!\left[\bigcup_{v\in S}B_v\right].
\]
Thus one circuit is equivalent to nonsingularity.
The count includes both choices at degree-two quotient vertices, whose blocks are empty.
If every selected degree is two, the quotient support is a cycle, all blocks are empty, and every one of its $2^{|V_+(F)|}$ cell choices produces one circuit.
This also agrees with the formula.
\end{proof}

The external identity used here is Theorem~4 of Lorenzo Traldi,
\href{https://arxiv.org/pdf/0903.4405}{\emph{Binary nullity, Euler circuits and interlace polynomials}},
European Journal of Combinatorics 32 (2011), 944--950.
It applies to undirected four-regular multigraphs, including loops.
Its transition rule deletes reference-following rows and columns and changes the diagonal for orientation-inconsistent transitions.

This representation alone does not give a half bound.
For example, an identity matrix has every principal submatrix nonsingular.
Even if the matrix is alternating, taking independent nonsingular $2$-by-$2$ blocks and permitting only whole-block selections gives the same phenomenon.
Any useful bound must therefore retain restrictions imposed by the original cubic graph and its cyclic connectivity.

There is a limited parity bound.
If $M_F$ has zero diagonal and at least one block has odd order, nonsingularity requires the total selected block order to be even.
Exactly half of all block subsets have that parity, so $\kappa(F)\le2^{|V_+(F)|-1}$.
Neither the zero-diagonal hypothesis nor the existence of an odd block is automatic.
